\section{Performance Evaluation}
\label{sec:performance}

\subsection{Optimized Implementations}
\label{sec:impl}

We evaluate three implementations of \TW{} sharing a common structure: a
portable reference written in Go, an \texttt{amd64} implementation, and an
\texttt{arm64} implementation. The two architecture-specific implementations
replace only the \Keccakp{} permutation and the inner absorb/squeeze
loops with hand-written assembly; the framing, domain separation, padding, and
tree bookkeeping remain in the shared Go code. The portable reference uses a
scalar permutation and a scalar eight-instance loop, and serves both as a
correctness oracle for the assembly cores and as the fallback on other
architectures.

The tree structure of \TW{} exposes two distinct permutation workloads, and
both architectures provide a separate core for each. The root duplex---which
absorbs the associated data, processes chunk~$0$, and absorbs the leaf tag
aggregation transcript---is inherently sequential, since each of its calls
consumes the previous call's output; it is driven by a \emph{single-duplex}
($\times 1$) core that permutes one 1600-bit state at a time, except for its
chunk-$0$ message phase, which is fused into the first batched leaf call as
described below. The leaf duplexes,
by contrast, operate on disjoint state and disjoint inputs (\Cref{sec:tw128}),
so the implementation batches eight complete leaf chunks and runs them in
lockstep through an \emph{eight-way} ($\times 8$) core. Because every leaf chunk
is exactly $\beta = 49\rho$ bytes, the eight instances step through an identical
sequence of $49$ full $\rho$-blocks (the first $48$ closed with
\textsc{msg\_more}, the last with \textsc{msg\_last}), so a single shared block
position drives all eight. The eight-way state is held lane-major---25 lanes
each holding the corresponding 64-bit word of all eight instances---so that the
batched cores can address the eight instances of a lane with a single
vector load.

Messages are decomposed into eight-chunk batches run on the $\times 8$ core,
and a leftover run of two to seven complete chunks is drained by a
narrower kernel: on \texttt{amd64} a register-resident kernel processes the
remainder in a single call, reading the active chunks in place, and on
\texttt{arm64} a $2$-wide kernel consumes it two chunks at a time. These
remainder kernels read each chunk as a sequential stream, so they prefetch as
well on the cold tail of a long message as on a cache-resident short one; on
the \texttt{amd64} host they outperform both repeated single-duplex calls and
padding the remainder to a full eight-chunk batch, at every remainder width.
Only a single leftover complete chunk, a partial trailing chunk, and the
count-aggregation tail fall back to the single-duplex ($\times 1$) core.

Chunk~$0$ itself runs on the batched kernels rather than serially, a technique
we refer to as \emph{lane-$0$ fusion}. The root duplex's chunk-$0$ message
phase has the same kernel-visible block schedule as a leaf chunk---the same
$49$ $\rho$-blocks with the same closing flags---and differs from a leaf only
in its initial state, to which the kernels are indifferent; chunk~$0$ is also
contiguous in the message with the first leaf chunks. The implementation
therefore seeds lane~$0$ of the first batched call with the root state (after
initialization and associated data processing) and runs the chunk-$0$ phase
alongside up to seven leaves: in the $\times 8$ core when the message has at
least seven complete leaf chunks, and in the \texttt{amd64} remainder kernel or
the \texttt{arm64} $2$-wide kernel below that. Lane~$0$'s output state is
written back to the root duplex, which then absorbs the leaf tags produced by
the same call; no serial pass over chunk~$0$ precedes the leaf work. On
\texttt{amd64} the assembly kernels handle only the
$48$ full \textsc{msg\_more} blocks of each chunk; the final \textsc{msg\_last}
block and the leaf tag extraction are performed in portable Go, shared with the
reference path, which keeps the byte-granular tail logic out of assembly. The
\texttt{arm64} kernels process the complete chunk, extracting the leaf tags and
storing the permutation states back in assembly---the writeback that allows the
root duplex to occupy lane~$0$.

\paragraph{Architecture-specific cores.}
On \texttt{amd64} the cores are built from AVX-512 (selected once at startup
from the \texttt{AVX512VL} CPUID feature bit) and AVX2, holding the lane-major
state in ZMM or YMM registers and permuting it with no cross-lane shuffles. The
steady-state $\times 8$ kernel transposes each rate block into the lane-major
layout with contiguous, prefetcher-friendly loads and stores, while the
register-resident kernel that drains a two-to-seven-chunk remainder reads its
active chunks in place by masked gather and scatter. On \texttt{arm64} the cores
use NEON with the Armv8.2 SHA-3 extension (\texttt{FEAT\_SHA3}, available on
Apple Silicon and Neoverse V1/V2 and N2), one instance per $64$-bit lane in the
single-duplex core and two instances per $128$-bit register in the $\times 8$
and $2$-wide kernels. Both batched paths store the permutation state back after
the closing block---the writeback that lets the root duplex occupy lane~$0$.
\Cref{app:kernels} gives the per-instruction detail, including the fused
permutation instructions and the remainder-kernel strategies on each
architecture.

\paragraph{Constant time.}
All three implementations are constant time with respect to the key, nonce, and
plaintext. The permutation and the XOR-based absorb/squeeze contain no
secret-dependent branches and no secret-dependent memory or vector indices. The
only indexed table is the round-constant array, addressed by the public round
number; the AVX-512 chunk transpose uses contiguous loads and stores, and the
gather and scatter that remain---for the partial rate lane and the
register-resident remainder kernel---are addressed by the fixed public chunk
strides under a mask set by the public chunk count, and the AVX2 remainder
kernel's redirection of inactive instances to its dummy block is likewise
selected by the public chunk count. None of these depend on secret data.

\subsection{Benchmark Methodology}
\label{sec:bench-method}

We measure on two hosts, one per architecture. The \texttt{amd64} host is a
Google Compute Engine \texttt{c4-standard-4} instance with an Intel Xeon Platinum
8581C (Emerald Rapids, $2.30$\,GHz base, AVX-512 including \texttt{AVX512VL},
\texttt{AVX512\_VBMI2}, \texttt{GFNI}, and \texttt{VAES}), two physical cores
with simultaneous multithreading enabled, $48$\,KiB L1d and $32$\,KiB L1i per
core, $2$\,MiB L2 per core, and a $260$\,MiB shared L3, running Debian~12 (kernel
6.1) under KVM. The guest does not expose a \texttt{cpufreq} interface, so the
frequency governor and turbo state are not under our control; we mitigate the
resulting variance statistically (below). The \texttt{arm64} host is an Apple M4
Pro ($10$ performance cores and $4$ efficiency cores; per performance core,
$128$\,KiB L1i, $64$\,KiB L1d, and a $16$\,MiB shared L2) running macOS~26.5.1
(Darwin~25.5), with \texttt{FEAT\_SHA3} present. macOS exposes no userspace
frequency control.

Measurements are collected by a standalone harness (\texttt{cmd/cpb}) that locks
its goroutine to an OS thread and, for each algorithm, operation, and message
length, first calibrates an iteration count that fills at least $100$\,ms, then
records $100$ samples, from which it reports the median cycles per byte and
median wall-clock throughput together with the interquartile range as a
dispersion measure (\Cref{tab:dispersion}). Both metrics are read from the same timed loop, so the
cycles-per-byte and throughput figures for a given measurement come from one
interleaved run rather than two separately scheduled ones. Each sample reuses the
same source and destination buffers, so the reported figures are warm-cache,
steady-state rates. The single-threaded measurement leaves the SMT sibling of the
\texttt{amd64} host idle. We report a sweep of seven message lengths from $1$\,B
to $16$\,MiB.

The cycle counters differ by architecture, and this affects how the
cycles-per-byte figures may be read. On \texttt{amd64} the harness reads
\texttt{RDTSC}; on this part the time-stamp counter is invariant and ticks at the
$2.30$\,GHz reference frequency, so the reported figures are \emph{reference}
cycles per byte and are unaffected by turbo (and convert directly to wall-clock
time). On \texttt{arm64} the harness reads \texttt{CNTVCT\_EL0} and scales it by
the ratio of the peak core frequency to \texttt{CNTFRQ\_EL0}; since Apple exposes
no frequency API, the peak is auto-detected by timing a dependent integer-add
chain and taking the top observed P-state, yielding \emph{core} cycles per byte.
Consequently the cycles-per-byte figures are comparable within an architecture
and against the same-host AES-128-GCM baseline, but not directly across
architectures; absolute throughput (\Cref{tab:throughput}) is the cross-platform
metric.

The baseline is Go's AES-128-GCM implementation
(\texttt{crypto/\allowbreak aes} with \texttt{crypto/\allowbreak cipher}).
It uses AES-NI and \texttt{PCLMULQDQ} on \texttt{amd64}, and the AES
and \texttt{PMULL} instructions on \texttt{arm64}.
Each timed AES-128-GCM call constructs its cipher, so its measurement includes
the key schedule and the GHASH key-power precomputation; this mirrors the
\TW{} measurement, whose one-shot seal and open perform the root duplex
initialization (one permutation) inside the timed call. Both directions are
measured for each scheme; the two directions track each other to within a few
percent across the sweep, with a worst-case divergence of $8\%$ (\TW{} on
\texttt{arm64} at $64$\,B).

\subsection{Throughput}
\label{sec:throughput}

\Cref{tab:cpb} reports cycles per byte across the message-length sweep and
\Cref{tab:throughput} the corresponding wall-clock throughput in the bulk
regime. The dominant feature of the \TW{} rows is the descent produced by the
tree-parallel leaf core: the per-byte cost falls as a message lengthens and its
chunks occupy more SIMD lanes. A message of at most one chunk has no complete
leaf chunk for chunk~$0$ to share a kernel with, so it runs entirely on the
single-duplex core, and the $8$\,KiB column sits at that core's rate. Past one
chunk, lane-$0$ fusion (\Cref{sec:impl}) keeps every complete chunk on the
batched kernels, and the descent tracks lane occupancy: at $32$\,KiB, chunk~$0$
and three complete leaves run four-wide in the \texttt{amd64} remainder kernels
and as two full pairs on \texttt{arm64}, and $64$\,KiB is the first length
whose eight complete chunks ($8\beta \approx 64$\,KiB) fill the eight-way
($\times 8$) core. On the \texttt{amd64} AVX-512 path \TW{} encryption falls
$1.89 \to 0.78 \to 0.40$ across the $8$, $32$, and $64$\,KiB columns and is
already at its bulk rate at $64$\,KiB, holding at $0.39$--$0.40$ through
$16$\,MiB. The AVX2 path follows the same shape at lower throughput ($2.20 \to
1.10$ over the same span, thereafter between $1.02$ and $1.08$); it runs the
eight instances as two groups of four rather than in a single $8$-wide
register, leaving the AVX-512 core about $2.7\times$ faster in bulk. The
\texttt{arm64} curve flattens earliest ($1.95 \to 0.91 \to 0.88 \to 0.88 \to
0.89$ cycles per byte from $8$\,KiB to $16$\,MiB): the eight-way core itself
processes its batch as four two-instance pairs, so the fused $2$-wide kernel
runs at essentially the same per-chunk rate, and a $32$\,KiB message---chunk~$0$
and three leaves as two full pairs---is already within $3\%$ of the $16$\,MiB
floor. In wall-clock terms \TW{} reaches $5.77$\,GB/s on the \texttt{amd64}
AVX-512 host and $5.04$\,GB/s on the M4 Pro for $16$\,MiB messages.

\begin{table}
\centering
\caption{Median cycles per byte ($100$ samples). The \texttt{amd64} figures are
\texttt{RDTSC} reference cycles at the $2.30$\,GHz invariant time-stamp counter;
the \texttt{arm64} figures are core cycles at the auto-detected peak performance-core
frequency. Cycle counts are not directly comparable across architectures
(\Cref{sec:bench-method}).}
\label{tab:cpb}
\begin{tabular}{l rrrrrrr}
\toprule
& \multicolumn{7}{c}{Message length} \\
\cmidrule(l){2-8}
& $1$\,B & $64$\,B & $8$\,KiB & $32$\,KiB & $64$\,KiB & $1$\,MiB & $16$\,MiB \\
\midrule
\multicolumn{8}{l}{\emph{Intel Xeon 8581C (Emerald Rapids), AVX-512}} \\
\quad \TW{} encrypt      & 590 & 9.27 & 1.89 & 0.78 & 0.40 & 0.39 & 0.40 \\
\quad \TW{} decrypt      & 615 & 9.72 & 1.90 & 0.75 & 0.40 & 0.39 & 0.39 \\
\quad AES-128-GCM seal   & 765 & 13.23 & 0.38 & 0.32 & 0.30 & 0.30 & 0.30 \\
\quad AES-128-GCM open   & 785 & 12.79 & 0.37 & 0.31 & 0.30 & 0.29 & 0.29 \\
\midrule
\multicolumn{8}{l}{\emph{Intel Xeon 8581C (Emerald Rapids), AVX2}} \\
\quad \TW{} encrypt      & 656 & 10.31 & 2.20 & 1.10 & 1.02 & 1.07 & 1.08 \\
\quad \TW{} decrypt      & 683 & 10.70 & 2.21 & 1.11 & 1.08 & 1.07 & 1.08 \\
\midrule
\multicolumn{8}{l}{\emph{Apple M4 Pro, NEON\,+\,\texttt{FEAT\_SHA3}}} \\
\quad \TW{} encrypt      & 620 & 9.77 & 1.95 & 0.91 & 0.88 & 0.88 & 0.89 \\
\quad \TW{} decrypt      & 667 & 10.59 & 2.00 & 0.89 & 0.87 & 0.87 & 0.87 \\
\quad AES-128-GCM seal   & 1044 & 17.28 & 0.54 & 0.45 & 0.44 & 0.43 & 0.44 \\
\quad AES-128-GCM open   & 1055 & 17.31 & 0.52 & 0.43 & 0.42 & 0.41 & 0.42 \\
\bottomrule
\end{tabular}
\end{table}

\begin{table}
\centering
\caption{Measured wall-clock throughput (GB/s, $\mathrm{GB} = 10^9$ bytes) from
the same \texttt{cmd/cpb} run as \Cref{tab:cpb}, encryption direction; the
decryption direction tracks it. The \texttt{amd64} AES-128-GCM row is the AVX-512
host and is unaffected by the \TW{} core selection.}
\label{tab:throughput}
\begin{tabular}{l rrrrr}
\toprule
& $8$\,KiB & $32$\,KiB & $64$\,KiB & $1$\,MiB & $16$\,MiB \\
\midrule
\multicolumn{6}{l}{\emph{Intel Xeon 8581C (Emerald Rapids)}} \\
\quad \TW{} (AVX-512)  & 1.22 & 2.95 & 5.78 & 5.94 & 5.77 \\
\quad \TW{} (AVX2)     & 1.05 & 2.09 & 2.25 & 2.15 & 2.14 \\
\quad AES-128-GCM      & 6.11 & 7.27 & 7.57 & 7.68 & 7.63 \\
\midrule
\multicolumn{6}{l}{\emph{Apple M4 Pro}} \\
\quad \TW{} (NEON+SHA3) & 2.30 & 4.92 & 5.06 & 5.10 & 5.04 \\
\quad AES-128-GCM       & 8.25 & 9.89 & 10.18 & 10.33 & 10.09 \\
\bottomrule
\end{tabular}
\end{table}

\paragraph{Dispersion.}
Each cell of \Cref{tab:cpb,tab:throughput} is the median of $100$ samples;
\Cref{tab:dispersion} summarizes their spread as the interquartile range
(IQR) relative to the median. The medians are stable: across the sweep the
relative IQR of the \TW{} cycles-per-byte measurements is at most $1.6\%$
on the \texttt{amd64} AVX-512 path, $2.1\%$ on AVX2, and $1.0\%$ on the M4
Pro, with per-path medians of $1.2\%$, $1.4\%$, and $0.6\%$. The full
min--max range is wider on the \texttt{amd64} guest---up to $16.5\%$
(AVX-512) and $20.2\%$ (AVX2) at isolated lengths---because its turbo and
governor states are not under our control (\Cref{sec:bench-method}); these
are tail excursions that the IQR excludes, and the M4 Pro, whose frequency
is steadier, keeps even its full range within $4.7\%$. We therefore report
medians throughout and treat the IQR as the dispersion measure.

\begin{table}
\centering
\caption{Cycles-per-byte measurement dispersion for \TW{} across the
seven-length sweep (seal and open, $100$ samples per point), as the
interquartile range relative to the median. The final column is the widest
single-point min--max range; its size on the \texttt{amd64} guest reflects
the turbo and governor states we cannot pin (\Cref{sec:bench-method}) and
is excluded from the IQR as a tail effect.}
\label{tab:dispersion}
\begin{tabular}{l rrr}
\toprule
& \multicolumn{2}{c}{Relative IQR} & Max \\
\cmidrule(lr){2-3}
Path & median & max & range \\
\midrule
Xeon 8581C, AVX-512 & $1.2\%$ & $1.6\%$ & $16.5\%$ \\
Xeon 8581C, AVX2    & $1.4\%$ & $2.1\%$ & $20.2\%$ \\
Apple M4 Pro        & $0.6\%$ & $1.0\%$ & $\phantom{0}4.7\%$ \\
\bottomrule
\end{tabular}
\end{table}

\subsection{Latency}
\label{sec:latency}

For short messages the cost is set by the serial root duplex and is independent
of the parallel leaf machinery. A message of at most $\rho = 167$ bytes occupies
a single rate block and never enters leaf mode, so it costs exactly two
permutations---one to initialize the root duplex and one to close the message
block and extract the tag. The $1$\,B and $64$\,B columns of \Cref{tab:cpb} bear
this out: per-message latency is nearly flat across them at roughly $590$
reference cycles on \texttt{amd64} ($590$ at $1$\,B versus $9.27 \times 64
\approx 593$ at $64$\,B) and $620$ core cycles on \texttt{arm64} ($620$ versus
$9.77 \times 64 \approx 625$), i.e.\ about $295$ and $310$ cycles per
\Keccakp{} respectively. Because such messages bypass the leaf
core entirely, leaf-level parallelism imposes no latency penalty on them.

The latency/throughput trade-off appears instead in the intermediate regime,
which lane-$0$ fusion has narrowed to roughly the one-to-eight-chunk band. The
worst per-byte cost in the sweep is the single-chunk regime at $8$\,KiB
($1.9$--$2.2$ cycles per byte across the three paths), where the message is the
chunk-$0$ phase alone on the single-duplex core. Once the message has complete
leaf chunks for chunk~$0$ to share a kernel with, the cost falls quickly: by
$32$\,KiB \TW{} is within a factor of two of its AVX-512 floor and within $3\%$
of its \texttt{arm64} floor, and the floor itself is reached once the eight-way
core fills at $64$\,KiB. The root duplex also bounds the tail of a long
message: the final tag cannot be produced until chunk~$0$ has been processed and
every leaf tag has been absorbed into the aggregation transcript. That absorption
is cheap---each rate block of the root absorbs five $32$-byte leaf tags, so a
$16$\,MiB message ($\approx 2000$ leaves) adds on the order of $400$ root
permutations against roughly $10^5$ permutations of leaf work, under $0.5\%$
overhead---so the serial root does not materially cap bulk throughput.

\subsection{Comparison with Other AEADs}
\label{sec:compare-perf}

We compare against the one baseline measured under the same harness on the same
hosts, the Go standard library's hardware-accelerated AES-128-GCM
(\Cref{tab:cpb,tab:throughput}). The harness times \TW{} and AES-128-GCM
back-to-back at each length and reads both metrics from the same loop, so the
cycles-per-byte and throughput comparisons hold under identical conditions and
agree. In the bulk regime AES-128-GCM is faster on both architectures, as
expected for a primitive with dedicated AES and carryless-multiply instructions:
at $16$\,MiB it runs at $0.30$ cycles per byte on \texttt{amd64} and $0.44$ on
\texttt{arm64}, against $0.40$ and $0.89$ for \TW{}---about $1.3\times$ and
$2\times$ faster. The wall-clock throughputs show the same margins: \TW{} reaches
$5.77$ and $5.04$\,GB/s against AES-128-GCM's $7.63$ and $10.09$\,GB/s. The wider
\texttt{arm64} gap reflects the strength of Apple's AES and \texttt{PMULL} units.
\TW{} attains these rates with a permutation that has no dedicated hardware
support, relying only on the SHA-3 extension on \texttt{arm64} and on general SIMD
on \texttt{amd64}.

The ordering reverses for very short messages, where AES-128-GCM's key schedule
and GHASH key-power precomputation dominate: at $1$\,B AES-128-GCM costs $765$
and $1044$ cycles on the two hosts against \TW{}'s $590$ and $620$. This
comparison is setup-bound rather than representative of steady state, since each
timed call rebuilds its cipher; it nonetheless reflects the per-message cost an
application pays when it does not amortize key setup across messages. In the
sub-bulk regime AES-128-GCM is several times faster at one chunk ($0.38$ versus
$1.89$ cycles per byte at $8$\,KiB on \texttt{amd64}), but the gap narrows
quickly once the fused kernels engage: about $2.4\times$ at $32$\,KiB, and the
bulk ratio from $64$\,KiB on. \TW{}'s intended value relative to AES-128-GCM
lies in the tree structure and committing security properties analyzed in
\Cref{sec:main-results,sec:cmtds3} rather than in raw speed.
